\section{Parametric description of lines}

The parametric description is one of the most natural ways to describe a line. It uses the idea that a line can be traced by a moving point.

We start with a fixed point and a fixed direction. Then we allow a parameter to vary. As the parameter changes, the point moves along the line.

\subsection*{Vector parametric form}

Let \(A\) be a point and let \(v\ne0\) be a direction vector. The line through \(A\) in direction \(v\) is
\[
P=A+t v,
\qquad t\in\mathbb{R}.
\]
This is a compact formula, but it contains three different roles:
\begin{itemize}
\item \(A\) is fixed and anchors the line,
\item \(v\) is fixed and gives the direction,
\item \(t\) varies and moves the point along the line.
\end{itemize}

If we do not distinguish these roles, the formula becomes only a string of symbols. If we do distinguish them, the formula becomes a description of motion.

\begin{examplebox}
Consider the line
\[
P=(2,1)+t(3,4).
\]
For \(t=0\), we get
\[
P=(2,1).
\]
For \(t=1\), we get
\[
P=(5,5).
\]
For \(t=-1\), we get
\[
P=(-1,-3).
\]
The same formula produces different points on the same line. The parameter controls how far we move in the direction \((3,4)\).
\end{examplebox}

The parameter is not something mysterious. It is a coordinate along the line, measured in copies of the direction vector.

\subsection*{Coordinate parametric form in the plane}

If
\[
A=(x_0,y_0),
\qquad
v=(a,b),
\]
then
\[
P=A+t v
\]
means
\[
(x,y)=(x_0,y_0)+t(a,b).
\]
Therefore
\[
x=x_0+ta,
\qquad
y=y_0+tb.
\]
These are the parametric equations of the line in the plane.

\begin{examplebox}
Let
\[
A=(-1,3),
\qquad
v=(2,-5).
\]
Then the line is
\[
(x,y)=(-1,3)+t(2,-5),
\]
or equivalently
\[
x=-1+2t,
\qquad
y=3-5t.
\]
For \(t=2\), the point is
\[
x=-1+4=3,
\qquad
y=3-10=-7.
\]
Thus \((3,-7)\) lies on the line.
\end{examplebox}

The coordinate equations should be read together. The same value of \(t\) must be used in both equations. We are not choosing one parameter for \(x\) and another parameter for \(y\). One moving point has one parameter.

\subsection*{Checking whether a point lies on a parametric line}

To check whether a point belongs to a parametric line, we ask whether there exists a value of the parameter that produces the point.

\begin{examplebox}
Consider the line
\[
x=1+3t,
\qquad
y=2-t.
\]
Does the point \((7,0)\) lie on the line?

From the first equation,
\[
7=1+3t,
\]
so
\[
t=2.
\]
Using the same value in the second equation gives
\[
y=2-2=0.
\]
Therefore \((7,0)\) lies on the line.

Now consider the point \((7,1)\). The first coordinate again forces \(t=2\), but then the second coordinate would be \(0\), not \(1\). Therefore \((7,1)\) does not lie on the line.
\end{examplebox}

This example shows why the parameter must be shared. A point lies on the line only if all coordinate equations agree on the same parameter value.

\subsection*{Different parametrizations of the same line}

A line does not have a unique parametric equation. We may choose a different point on the same line, or a different non-zero vector parallel to the same direction.

For example,
\[
P=(2,1)+t(3,4)
\]
and
\[
P=(5,5)+s(6,8)
\]
describe the same line. The point \((5,5)\) lies on the first line, and the vector \((6,8)\) is twice the direction vector \((3,4)\). The second description moves along the same line, but uses a different starting point and a different scale for the parameter.

\begin{remarkbox}
A parametric equation describes a line, but the parametrization itself is not unique. The geometric line remains the same even if we choose another point on it or replace the direction vector by a non-zero scalar multiple.
\end{remarkbox}

\subsection*{Why parametric form is useful}

Parametric form is especially useful when the natural data are a point and a direction, or two points. It is also the most natural form for lines in space, where slope-intercept equations are no longer available.

\begin{summarybox}
When reading a parametric line, ask:
\begin{itemize}
\item What point is fixed on the line?
\item What vector gives the direction?
\item What does the parameter do geometrically?
\item Does a given point arise from one common parameter value?
\item Could a different point or parallel direction vector describe the same line?
\end{itemize}
\end{summarybox}

The parametric form is not merely a formula. It is a way to watch a point move.